%====================================
% Environment & debugging
% \linenumbers % Show line numbers
%====================================
\DTMlangsetup[en-GB]{ord=raise,monthyearsep={,\space}} % date format


%====================================
% Structure & Layout
%====================================
% Better page layout for A4 paper, see memoir manual.
\settrimmedsize{297mm}{210mm}{*}
\setlength{\trimtop}{0pt} 
\setlength{\trimedge}{\stockwidth} 
\addtolength{\trimedge}{-\paperwidth} 
\settypeblocksize{634pt}{448.13pt}{*} 
\setulmargins{4cm}{*}{*} 
\setlrmargins{*}{*}{1.5} 
\setmarginnotes{17pt}{51pt}{\onelineskip} 
\setheadfoot{2\onelineskip}{2\onelineskip} 
\setheaderspaces{*}{2\onelineskip}{*} 
\checkandfixthelayout

% Page layout
\makepagestyle{myvf} 
\makeoddfoot{myvf}{}{\thepage}{} 
\makeevenfoot{myvf}{}{\thepage}{} 
\makeheadrule{myvf}{\textwidth}{\normalrulethickness} 
\makeevenhead{myvf}{\small\textsc{\leftmark}}{}{} 
\makeoddhead{myvf}{}{}{\small\textsc{\rightmark}}
\pagestyle{myvf}




%====================================
% Font & Spacing
%====================================
\sloppy % This will make LaTeX adjust spacing to avoid overfull lines

\setSingleSpace{1.6}
\SingleSpacing
\frenchspacing


% setting heading line spacing
\makeatletter
\renewcommand\section{\@startsection{section}{1}{\z@}%
  {-3.5ex \@plus -1ex \@minus -.2ex}%
  {2.3ex \@plus .2ex}%
  {\normalfont\Large\bfseries\renewcommand{\baselinestretch}{1}\selectfont}}

\renewcommand\subsection{\@startsection{subsection}{2}{\z@}%
  {-3.25ex \@plus -1ex \@minus -.2ex}%
  {1.5ex \@plus .2ex}%
  {\normalfont\large\bfseries\renewcommand{\baselinestretch}{1}\selectfont}}

\renewcommand\subsubsection{\@startsection{subsubsection}{3}{\z@}%
  {-3.25ex \@plus -1ex \@minus -.2ex}%
  {1.5ex \@plus .2ex}%
  {\normalfont\normalsize\bfseries\renewcommand{\baselinestretch}{1}\selectfont}}

\makeatother





% epigraph
\epigraphfontsize{\small\itshape}
\setlength\epigraphwidth{8cm}
\setlength\epigraphrule{0pt}


% Reduce widows (the last line of a paragraph at the start of a page) and orphans (the first line of paragraph at the end of a page)
\widowpenalty=1000
\clubpenalty=1000


%====================================
% Chapters
%====================================
\makeatletter 
\newlength\dlf@normtxtw % do we need this???
\setlength\dlf@normtxtw{\textwidth} % do we need this???

% Chapter boxed number
\newsavebox{\feline@chapter} 
\newcommand\feline@chapter@marker[1][4cm]{%
	\sbox\feline@chapter{% 
		\resizebox{!}{#1}{\fboxsep=1pt%
			\colorbox{gray}{\color{white}\thechapter}% 
		}}%
		\rotatebox{90}{% Verticle Chaper
			\resizebox{%
				\heightof{\usebox{\feline@chapter}}+\depthof{\usebox{\feline@chapter}}}% 
			{!}{\scshape\so\@chapapp}}\quad%
		\raisebox{\depthof{\usebox{\feline@chapter}}}{\usebox{\feline@chapter}} % boxed number
} 

\newcommand\feline@chm[1][4cm]{%
	\sbox\feline@chapter{\feline@chapter@marker[#1]}% 
	\makebox[0pt][c]{% aka \rlap
		\makebox[1cm][r]{\usebox\feline@chapter}%
	}}


\makechapterstyle{daleifmodif}{
	\renewcommand\chapnamefont{\normalfont\Large\scshape\raggedleft\so\renewcommand{\baselinestretch}{1}\selectfont} 
	\renewcommand\chaptitlefont{\normalfont\Large\bfseries\scshape\renewcommand{\baselinestretch}{1}\selectfont} 
	\renewcommand\chapternamenum{} 
	\renewcommand\printchaptername{} 
	\renewcommand\printchapternum{\null\hfill\feline@chm[2.5cm]\par} % box height
	\renewcommand\afterchapternum{\par\vskip\midchapskip} 
	\renewcommand\printchaptertitle[1]{\color{gray}\chaptitlefont\raggedleft ##1\par}
}

% \makechapterstyle{daleifmodif}{
% 	\renewcommand\chapnamefont{\normalfont\Large\scshape\raggedleft\so} 
% 	\renewcommand\chaptitlefont{\normalfont\Large\bfseries\scshape} 
% 	\renewcommand\chapternamenum{} \renewcommand\printchaptername{} 
% 	\renewcommand\printchapternum{\null\hfill\feline@chm[2.5cm]\par} % box height
% 	\renewcommand\afterchapternum{\par\vskip\midchapskip} 
% 	\renewcommand\printchaptertitle[1]{\color{gray}\chaptitlefont\raggedleft ##1\par}
% } 
\makeatother 
% end modifying package internals

\chapterstyle{daleifmodif}




%====================================
% Sections
%====================================
% Sets numbering division level
\setsecnumdepth{subsection} 
\maxsecnumdepth{subsubsection}


%====================================
% Tables & Figures
%====================================
% Figure/table as subfloat
\newsubfloat{figure}
\newsubfloat{table}


% caption: bold type, italic label
\captiontitlefont{\itshape} 
\captionnamefont{\bfseries}

% set table and figure white space above and line spacing
\AtBeginEnvironment{figure}{\vspace{10pt}}
\AtBeginEnvironment{figure}{\setSingleSpace{1}\selectfont}

\AtBeginEnvironment{table}{\vspace{18pt}}
\AtBeginEnvironment{table}{\setSingleSpace{1}\selectfont}
\AtEndEnvironment{table}{\vspace{10pt}}


% 
\AtBeginEnvironment{longtable}{%
  \begingroup
  \renewcommand{\baselinestretch}{1}\selectfont
}
\AfterEndEnvironment{longtable}{\endgroup}

% table cell padding
\renewcommand{\arraystretch}{1.50}





%====================================
% Flowcharts
%====================================
% Define a custom color using a hex value
%\definecolor{customGreen}{HTML}{006400}
%\usetikzlibrary{arrows.meta,decorations.markings}  % for arrow heads
\usetikzlibrary{arrows.meta,decorations.markings, fit, positioning}

\tikzstyle{start} = [ellipse,  
    text width=1.75cm,
    minimum height=1.25cm, 
    line width=1pt, 
    text centered, 
    draw=green!40!black!40, 
    fill=green!40!black!30, 
    font=\scriptsize
]

\tikzstyle{stop} = [
    ellipse,  
    text width=1.75cm,
    minimum height=1.25cm, 
    line width=1pt, 
    text centered, 
    draw=red!40, 
    fill=red!25, 
    font=\scriptsize
]


\tikzstyle{process} = [
    rectangle, 
    minimum height=1.5cm, 
    line width=1pt, 
    text centered, 
    text width=3.75cm, 
    draw=blue!30, 
    fill=blue!15, 
    align=center,
    font=\linespread{0.9}\scriptsize
]


\tikzstyle{decision} = [
    diamond, 
    aspect=1.75, 
    minimum height=2.75cm, 
    line width=1pt, 
    text centered, 
    text width=2.75cm, 
    draw=yellow!60, 
    fill=yellow!40, 
    font=\linespread{0.9}\scriptsize
]


\tikzstyle{io} = [
    draw=none,
    inner sep=0pt,
    outer sep=0pt,
    minimum width=4.00cm,
    minimum height=1.5cm,
    font=\linespread{0.9}\scriptsize\selectfont,
    align=center,
    path picture={
        \filldraw[draw=orange!50, fill=orange!30]
        ([xshift=0.50cm]path picture bounding box.north west) --
        ([xshift=0.00cm]path picture bounding box.north east) --
        ([xshift=-0.50cm]path picture bounding box.south east) --
        ([xshift=0.00cm]path picture bounding box.south west) -- cycle;
    }
]


\tikzstyle{arrow} = [->,>=stealth, font=\scriptsize, line width=1pt]

